\subsection*{Example Code}

\begin{lstlisting}[language=croqtile,caption={Example Code}]
__co__ auto rgb2gray(f32 [N, 3, H, W] input) {
  f32 [N, H, W] out;
  parallel q by 6 {
    with index={n, h, w} in [N, H, W]/{#q, 16, 512} {
      foreach n, h, w {
        // _ means no tiling
        input_L1_A = dma.copy input.chunkat(q#n, _, h, w) => local;
        dims : input_L1_A.span[(0), (2), (3)];
        local f32 [dims] out_L1;
        call rgb2gray_kernel_fp32_fp32(input_L1_A.data, out_L1, |out_L1|, 1);
        dma.copy out_L1 => out.chunkat(q#n, h, w);
      }
    }
  }
  return out;
}
\end{lstlisting}

The input shape is \co@[N, 3, H, W]@, with the three color channels in the second dimension. The underscore in \co@chunkat(q#n, _, h, w)@ keeps those channels together while selecting one image and a spatial region. Removing that channel dimension from the loaded shape gives the shape of \co@out_L1@; the helper writes one grayscale value for each spatial position.

The composition \co@q # n@ assigns a consecutive group of images to each of the six parallel instances. The bounds assume that N divides by 6, H by 16, and W by 512 for this regular partition. Within those assumptions, each input tile contains all three channels for a 16 by 512 region, and the final copy writes its grayscale result to the matching output tile. Other image sizes need an explicit boundary policy.
